\section{Experimental protocol}
\label{sec:protocol}

\paragraph{Datasets.}
We use five public, human-labelled English datasets that represent common decision tasks
(\hyperref[tab:datasets]{Table~\ref*{tab:datasets}}): intent classification for assistants, CLINC150 in-scope \citep{larson2019evaluation} and
HWU64 \citep{liu2019benchmarking}; financial sentiment, FPB \citep{malo2014good}; hate and offensive speech,
HateXplain \citep{mathew2021hatexplain}; and question type, TREC-6 \citep{li2002learning}. CLINC150, HateXplain and
TREC-6 use their official splits. FPB has no official split; we use a public 64/16/20 split of the 50\%-agreement subset,
following \citet{araci2019finbert}, whose 970-example test set is the one used by \citet{deusser2026jev}. For HWU64 and TREC-6, which have no official development
split, the development set is a stratified sample of the training set that follows the training class distribution.
Train and development texts that duplicate a test text after normalisation (lower-casing and removing punctuation),
and texts that occur with conflicting labels, are removed; the test sets are never changed. All raw files and processed
splits are pinned by SHA-256 hashes.

\begin{table}[t]
  \centering
  \caption{\textbf{The five datasets span intent, sentiment, moderation and question-type decisions with 3 to 150
  classes.} CLINC150 and HWU64 use fixed label groups as candidate sets ($K$ = 18--22 options); the other three use all
  labels. Train/dev/test sizes after cleaning.}
  \label{tab:datasets}
  \scriptsize
  \begin{tabular}{llrrrrr}
\toprule
Dataset & Domain & Classes & Candidates $K$ & Train & Dev & Test \\
\midrule
CLINC150 & intent (assistant) & 150 & 18--19 & 14,939 & 2,994 & 4,500 \\
HWU64 & intent (home assistant) & 64 & 21--22 & 7,871 & 1,076 & 1,076 \\
FPB & financial sentiment & 3 & 3 & 3,090 & 775 & 970 \\
HateXplain & hate / offensive speech & 3 & 3 & 15,330 & 1,921 & 1,924 \\
TREC-6 & question type & 6 & 6 & 4,869 & 498 & 500 \\
\bottomrule
\end{tabular}

\end{table}

\paragraph{Label texts and candidate sets.}
Every label has a short name and a one-sentence description, written once for all systems and frozen before any model
was run. For datasets with at most 20 labels, the candidate set contains all labels. For CLINC150 and HWU64 we
partition the labels once into fixed groups (8 groups of 18--19 and 3 groups of 21--22) by balanced $k$-means on
embeddings of the label texts, and the candidate set of an example is the group of its gold label. We chose fixed groups
because a leakage audit showed that candidate sets built around the gold label (the gold label plus its nearest
neighbours) reveal the gold label from the option list alone: a classifier that never sees the text identifies it for
97.6\% of CLINC150 and 90.2\% of HWU64 examples, against 5\% chance. With fixed groups, every member of a group receives
the same set of options. The order of the options is shuffled per example with a fixed seed, with the gold position
counterbalanced across examples, and all systems read the same candidate file.

\paragraph{Systems and inputs.}
\hyperref[tab:systems]{Table~\ref*{tab:systems}} lists the systems. Jev is queried through its hosted API, with the label names as option keys
and the label descriptions as criteria; the response's version field was \texttt{jev-1.13.0} for every call (16{,}234 in the full run and 7{,}500 in
three repeats); 465 of the 16{,}234 full-run calls (464 rate-limit responses and one timeout) succeeded after one
retry, and no call failed. Jev's API returns its probabilities rounded to two decimals; we use them as returned and renormalise over the candidates. Each open TDM and its base model receive \emph{token-identical} prompts
in the TDM's own training format: the Tev1 format for Tev1-4B and Qwen3.5-4B, and the Nimble format for Nimble-9B and
Qwen3.5-9B (checked on all 16{,}234 development and test prompts). Thinking mode is disabled. The probability vector is
the softmax ($T=1$) of the candidate letters' logits at the first answer position, computed in fp32 from the final
hidden state. The two 9B models run in 8-bit precision so that they fit on a 12\,GB GPU, and the Nimble adapter is
checked to be fully loaded before every run; the two 4B models run in bf16. Local scoring was deterministic in our check:
scoring 50 CLINC150 development examples twice gave identical probabilities for all four local models. Tev1 was trained mostly on decisions with at most
eight options, so the 18--22 options of CLINC150 and HWU64 lie outside its training distribution, and the licence of its
weights was still being finalised at the time of writing. Nimble-9B is the \texttt{original-2676} release, trained on
2{,}676 synthetic examples; we do not use the later release whose training data include public TREC questions
(\hyperref[app:runs]{Appendix~\ref*{app:runs}}).

\begin{table}[t]
  \centering
  \caption{\textbf{Each open TDM is paired with its own backbone under identical prompts, readout and precision.}
  Tev1-4B with Qwen3.5-4B (Tev1 prompt format) and Nimble-9B with Qwen3.5-9B (Nimble prompt format, identical 8-bit base
  weights). Revisions are Hugging Face commit hashes. Readout: ``letters'' is the softmax ($T=1$) over the logits of the
  candidate letters at the first answer position; the trained classifiers use a softmax over each example's candidates;
  Jev returns probabilities directly.}
  \label{tab:systems}
  \scriptsize
  \setlength{\tabcolsep}{4pt}
  \begin{tabular}{lllll}
\toprule
System & Source & Revision & Precision & Readout \\
\midrule
\textsc{Typed decision models} & & & & \\
Jev & Cloudflare API (\texttt{typesafe/jev}) & \texttt{jev-1.13.0} & API & API probabilities \\
Tev1-4B & \texttt{togethercomputer/Tev1-4B-experimental} & \texttt{0b7becf0} & bf16 & letters, Tev1 format \\
Nimble-9B & \makecell[l]{Qwen3.5-9B + LoRA \\ \texttt{bespokelabs/Bespoke-Nimble-9B}} & \makecell[l]{\texttt{fe33f82b} \\ (\texttt{original-2676})} & 8-bit & letters, Nimble format \\
\midrule
\textsc{Base LLMs} & & & & \\
Qwen3.5-4B & \texttt{Qwen/Qwen3.5-4B} & \texttt{851bf6e8} & bf16 & as Tev1-4B \\
Qwen3.5-9B & \texttt{Qwen/Qwen3.5-9B} & \texttt{c2022362} & 8-bit & as Nimble-9B \\
\midrule
\textsc{Trained classifiers} & & & & \\
DeBERTa-v3-large & fine-tuned on the full training set, 5 seeds & \texttt{64a8c8ea} & fp16 & softmax \\
TF-IDF + LR & word 1--2 grams, $C$ selected on dev & scikit-learn 1.9.1 & CPU & softmax \\
\bottomrule
\end{tabular}

\end{table}

\paragraph{Trained classifiers.}
DeBERTa-v3-large \citep{he2023debertav3} is fine-tuned on the full training set with AdamW, a learning rate selected
from $\{10^{-5}, 2\cdot10^{-5}\}$ on the development set, 10\% linear warm-up, weight decay 0.01, batch size 16, mixed
precision, maximum length 128 and 5 epochs. The best checkpoint is selected by development Macro-F1 within the candidate
sets, evaluated every 100 steps. We train five seeds; a run would count as failed if its development accuracy did not
exceed that of a majority-class predictor \citep{dodge2020finetuning,mosbach2021stability}, which did not happen. All analyses use the first
seed; the other seeds are reported in \hyperref[tab:seeds]{Table~\ref*{tab:seeds}}. The TF-IDF baseline is a logistic regression on word
uni- and bigrams with sublinear term frequency; $C$ is selected on the development set from
$\{0.1, 1, 10, 10^2, 10^3, 10^4\}$, and FPB and TREC-6 select the largest value; extending the grid to $10^5$ and
$10^6$ changes their development Macro-F1 by at most 0.2~points and their test Macro-F1 by at most 0.7~points
(exploratory check). Both classifiers' probabilities are
renormalised over each example's candidate set. Note an asymmetry: the trained classifiers use the development set both
for model selection and for recalibration, whereas the zero-shot systems use it only for recalibration and the implicit
prior.

\paragraph{Label budgets.}
For RQ4, both classifiers are also trained on stratified subsets of 100, 250, 500, 1{,}000 and 2{,}500 training
examples, and of 5{,}000 where the training set is larger (HateXplain, HWU64, CLINC150), with five subset seeds per
budget; the full training set is the last point of each curve, with all five DeBERTa-v3-large seeds. Class quotas are
proportional to the training frequencies, with at least one example per class where the budget allows it; at 100
labels, CLINC150 has examples for only 100 of its 150 classes (\hyperref[tab:shots]{Table~\ref*{tab:shots}}).
DeBERTa-v3-large uses the settings above with a fixed schedule: 20 epochs up to 500 labels and 5 epochs above, with at
least 500 optimizer steps, following the advice for few-sample fine-tuning
\citep{mosbach2021stability,zhang2021revisiting}. The learning rate, $C$ and the checkpoint are selected on the full
development set at every budget. A failed run would be restarted with one of three reserve seeds per dataset and
budget; none failed.

\paragraph{Task fine-tuning of the base model.}
Condition (b) of the decision-training comparison fine-tunes Qwen3.5-4B on the target task with QLoRA
\citep{hu2022lora,dettmers2023qlora} at 100, 500 and 2{,}500 labels, using the same subsets as above and three seeds. The
adapter has rank 16, $\alpha=32$ and dropout 0.05 on the same projection matrices as the Nimble-9B adapter; it is trained
on a 4-bit NF4 copy of the base with learning rate $10^{-4}$, effective batch size 16, the DeBERTa schedule and a
cross-entropy loss over the candidate letters at the answer position. Checkpoints are selected by development Macro-F1
every 100 steps. Every reported prediction comes from the kept adapter on the bf16 base, with the same prompt, assistant
prefix and letter readout as Qwen3.5-4B and Tev1-4B, so the three 4B conditions (a) frozen base, (b) task QLoRA and (c)
Tev1-4B differ only in their weights.

\paragraph{Statistical tests.}
The 35 confirmatory tests are P1 and P2 (Macro-F1 and ECE on five datasets each; 20 tests) and PA1--PA3 (five datasets
each; 15 tests). P1 and P2 use a paired bootstrap over test examples with $B=10{,}000$ replicates; both systems and both
metrics use the same resampled indices. For PA1--PA3 (\hyperref[eq:pa1]{Eqs.~\ref*{eq:pa1}}--\nobreak\ref{eq:pa3}), the estimand is the mean over a fixed set of 50 target priors
drawn once at $\alpha=1$ (centred on the training prior); it is therefore conditional on this set of priors, and the
intervals reflect only the finiteness of the test set. Each of $B=10{,}000$ replicates draws a class-stratified bootstrap
sample of the test set and builds all 50 shifted samples from it by resampling with replacement, so that Jev and
DeBERTa-v3-large are evaluated on identical samples. PA1 compares each shifted ECE with a no-shift control drawn by the
same resampler at the base sample's own class prior, because top-label ECE is biased upward on resampled data. PA2
compares risks, which are nearly unbiased under resampling, with the base sample itself. The point estimates of P1 and P2 are the differences on the full test set, those of
PA1--PA3 the means of the bootstrap replicates. We report percentile 95\% intervals,
which are not adjusted for multiplicity, and two-sided $p=\min\{1,\,2\min\{\Pr(\Delta\le0),\Pr(\Delta\ge0)\}\}$, adjusted with the
Holm procedure \citep{holm1979simple} over all 35 tests (also for PA3, whose hypothesis is directional). With
$B=10{,}000$, a raw $p$ of zero is reported as $p<0.0002$ $(=2/B)$, and its Holm-adjusted value as $p<0.007$
$(=35\cdot2/B)$. Quantities fitted on the development set (thresholds, implicit priors, temperatures) are held fixed
within the bootstrap, so the intervals do not include development-set uncertainty. The descriptive prevalence-shift
experiment uses 200 Dirichlet draws per $\alpha$ for all systems and all five arms, with and without replacement.

\paragraph{Protocol and runs.}
Hypotheses, tests, the comparison family, all configurations and the code were fixed and recorded as SHA-256 hashes
before the final runs, and every final run was verified against these hashes. One amendment followed three
code audits; it changed the analysis, the DeBERTa failed-run rule and the Jev error handling, but not the code used by
the local language models. It was made after the first local test predictions had been produced, but before the
DeBERTa and Jev runs and before any of the 35 tests or any shift analysis was computed on the final runs
(\hyperref[app:deviations]{Appendix~\ref*{app:deviations}}). Before the freeze, some pipeline checks were run on test
splits, mainly FPB's, including early versions of PA1--PA3 and of the shift experiment on FPB;
\hyperref[app:deviations]{Appendix~\ref*{app:deviations}} discloses them.
The local models and both trained classifiers were run on one NVIDIA RTX 4070 Ti (12\,GB); Jev, a hosted model, was
queried through its API for the main runs in about two hours, inside the protocol's 72-hour window (\hyperref[app:runs]{Appendix~\ref*{app:runs}}).
The label-budget study was specified with the rest of the protocol; its QLoRA settings, trainer and curve analysis were
fixed and hashed in a second freeze before any of its final test predictions, after a development-only pilot that left every
setting unchanged (\hyperref[app:deviations]{Appendix~\ref*{app:deviations}}). It reuses the Jev and zero-shot
predictions of the main study.
